\section{Cartes régionales du taux de privation par dimension}
\label{ann:F}

Les cartes ci-dessous complètent la carte d'incidence N-MODA de la
section~\ref{fig:carte_nmoda} en présentant, pour chacune des sept
dimensions du bien-être de l'enfant, le taux de privation régional aux
deux vagues (EHCVM~I à gauche, EHCVM~II à droite). L'échelle de couleur
est commune (0 à 100\,\%), ce qui permet de comparer directement les
dimensions entre elles et les deux périodes. Les régions les plus
sombres sont les plus privées.

% Chaque carte est pivotee a 90 degres (contenu non flottant, une par page)
% pour occuper la largeur d'une orientation paysage sans inserer de page
% blanche : l'introduction ci-dessus reste en portrait.
\newcommand{\cartedim}[3]{%
  \clearpage
  \vspace*{\fill}
  \begin{center}
  \rotatebox{90}{%
    \begin{minipage}{0.85\textheight}
      \centering
      \includegraphics[width=\linewidth]{figures/#1}\par\smallskip
      \captionof{figure}{Taux de privation : #2}\label{#3}\par
      {\footnotesize \textit{Source :} Calcul de l'auteur, EHCVM.}
    \end{minipage}}
  \end{center}
  \vspace*{\fill}
}

\cartedim{fig_carte_dim_assai.pdf}{assainissement}{fig:carte_dim_assai}
\cartedim{fig_carte_dim_eau.pdf}{eau}{fig:carte_dim_eau}
\cartedim{fig_carte_dim_logem.pdf}{logement}{fig:carte_dim_logem}
\cartedim{fig_carte_dim_nutri.pdf}{nutrition}{fig:carte_dim_nutri}
\cartedim{fig_carte_dim_sante.pdf}{santé}{fig:carte_dim_sante}
\cartedim{fig_carte_dim_protect.pdf}{protection}{fig:carte_dim_protect}
\cartedim{fig_carte_dim_educ.pdf}{éducation}{fig:carte_dim_educ}
